\documentclass[11pt,reqno]{amsart}
\usepackage[T1]{fontenc}
\usepackage{lmodern}
\usepackage{microtype}
\usepackage{amsmath,amssymb,mathtools}
\usepackage[colorlinks=true,linkcolor=blue,citecolor=blue,urlcolor=blue]{hyperref}

\newtheorem{theorem}{Theorem}[section]
\newtheorem{lemma}[theorem]{Lemma}
\newtheorem{corollary}[theorem]{Corollary}
\newtheorem{remark}[theorem]{Remark}

\newcommand{\li}{\operatorname{li}}
\newcommand{\lcm}{\operatorname{lcm}}
\newcommand{\e}{\mathrm e}
\newcommand{\cS}{\mathcal S}
\newcommand{\cN}{\mathcal N}
\newcommand{\cK}{\mathcal K}
\newcommand{\cJ}{\mathcal J}

% TODO(submission): author metadata.
\title[A loglog improvement of Vaughan's bound]{A logarithmic improvement of Vaughan's bound\\ for the Erd\H{o}s--Straus exceptional set}
\author{Anonymous}
\date{Draft, 6 September 2026}

\begin{document}

\begin{abstract}
Let $E(N)$ be the number of $n\le N$ for which $4/n$ is not a sum of three unit
fractions.  Vaughan (1970) proved $E(N)\ll N\exp\{-c(\log N)^{2/3}\}$.  We
show
\[
 E(N)\ll N\exp\{-c(\log N)^{2/3}(\log\log N)^{1/3}\}.
\]
The proof is Vaughan's large-sieve argument applied to the residue classes
modulo $k\ell$ (rather than modulo $\ell$ alone) that force solvability,
where $\ell$ is a prime and $k\le \delta\log N$ a ``multiplier'', after
splitting the sieved set into progressions modulo $\lcm\{k\}$.  The new
ingredient is a count of the forced classes, based on an elementary lattice
estimate, which is uniform in the progression.  Only Bombieri--Vinogradov,
Brun--Titchmarsh and Montgomery's large sieve are used.
\end{abstract}

\maketitle

\section{Introduction}

The Erd\H{o}s--Straus conjecture asserts that for every integer $n\ge2$ there
are positive integers $x,y,z$ with
\begin{equation}\label{ES}
 \frac4n=\frac1x+\frac1y+\frac1z .
\end{equation}
It has been verified for $n\le10^{18}$ (see \cite{PW}, \S6, and \cite{Salez}
for $10^{17}$), and Mordell's identities show that it holds unless
$n\equiv 1^2,11^2,13^2,17^2,19^2,23^2\pmod{840}$ \cite{Mordell}.  Let
\begin{gather*}
 E(N)=\#\{n\le N:\ \eqref{ES}\ \text{has no solution}\},\\
 E_{\mathrm{pr}}(N)=\#\{p\le N\ \text{prime}:\ \eqref{ES}\ \text{has no solution}\}.
\end{gather*}
Vaughan \cite{Vaughan} proved that $E(N)\ll N\exp\{-c(\log N)^{2/3}\}$ for
an absolute $c>0$; a self-contained account of his argument, uniform in the
numerator, is given by Pomerance and Weingartner \cite{PW}, whose
Theorem~1.3 retains this exponent for fixed numerators.  We prove:

\begin{theorem}\label{thm:main}
There is an absolute constant $c>0$ such that for all $N\ge3$
\[
 E_{\mathrm{pr}}(N)\ll N\exp\{-c(\log N)^{2/3}(\log\log N)^{1/3}\}.
\]
\end{theorem}

\begin{theorem}\label{thm:all}
The same bound holds for $E(N)$.
\end{theorem}

The exponent $2/3$ is unchanged; the gain is the factor $(\log\log N)^{1/3}$
inside the exponential.  We briefly recall where Vaughan's $2/3$ comes from,
because the improvement is located precisely there.  For a prime
$\ell\equiv3\pmod4$, every $n$ in one of roughly $\tau((\ell+1)^2/16)$
explicit residue classes modulo $\ell$ satisfies \eqref{ES}; on average over
$\ell$ this is $\asymp(\log\ell)^2$ classes, so
$\sum_{\ell\le X}\omega(\ell)/\ell\asymp(\log X)^2$.  Feeding these classes
into Montgomery's large sieve saves a factor $\exp\{c(\log X)^2\}$, and the
sieve level constraint forces $\log X\asymp(\log N)^{1/3}$.

The observation here is that for each ``multiplier'' $k\equiv1\pmod4$ the
number $(k\ell+1)/4$ also produces forced classes, now modulo $k\ell$
(Lemma~\ref{lem:identity}).  Restrict $n$ to a fixed reduced progression
modulo $L=\lcm\{k\le K,\ k\equiv1\ (4)\}$.  The compatible classes become
classes modulo $\ell$; the small coprime factorizations selected in
Lemma~\ref{lem:mass} give distinct classes, even for different $k$.
Their averaged contribution multiplies the
class mass by
\[
 h(K)=\sum_{\substack{k\le K\\ k\equiv1\,(4)}}\frac{\varphi(k)}{k^2}\asymp\log K .
\]
The price is that the large sieve is applied to an interval of length
$N/L$; we choose $L\le N^{1/2}$, ensured by taking $K$ a small constant
times $\log N$.  Hence the gain
$\log K\asymp\log\log N$ and, after re-optimising, the exponent
$(\log N)^{2/3}(\log\log N)^{1/3}$.  Since
\[
 \sum_{k\mid L}\frac{\varphi(k)}{k^2}\le\prod_{p\mid L}(1-1/p)^{-1}\ll\log\log(3L),
\]
no multiplier family with $L\le N^{1/2}$ improves the order of this
particular weighted multiplier sum beyond $\log\log N$.

The argument is short.  Sections~\ref{sec:identity}--\ref{sec:mass} construct
and count the forced classes; Section~\ref{sec:sieve} runs the large sieve;
Section~\ref{sec:all} deduces Theorem~\ref{thm:all} from
Theorem~\ref{thm:main} by a Rankin argument on the semigroup generated by
exceptional primes.  Throughout, $p$ and $\ell$ denote primes, $\tau$ is the
divisor function, $\varphi$ Euler's function, and all constants are
absolute unless indicated.

\section{Forced residue classes}\label{sec:identity}

\begin{lemma}[multiplier identity]\label{lem:identity}
Let $k,\ell$ be positive integers with $k\ell\equiv3\pmod 4$, and let
$A=(k\ell+1)/4=uvw$ be any factorization into positive integers.  If $n\ge1$
satisfies $nv\equiv-u\pmod{k\ell}$, then $s=(nv+u)/(k\ell)$ is a positive
integer and
\[
 \frac4n=\frac1{suw}+\frac1{nsvw}+\frac1{nuvw}.
\]
In particular \eqref{ES} is solvable for every $n$ in the residue class
$-uv^{-1}\pmod{k\ell}$ (note $(v,k\ell)=1$ since $(A,k\ell)=1$).
\end{lemma}

\begin{proof}
On the common denominator $nsuvw$ the three numerators sum to
$nv+u+s=sk\ell+s=s(k\ell+1)=4suvw$.
\end{proof}

For $k=1$ these are the classes used by Vaughan.  We shall use primes
$\ell\equiv3\pmod4$ and multipliers $k\equiv1\pmod4$.  The condition
$nv\equiv-u\pmod{k\ell}$ splits, when $(k,\ell)=1$, into
\begin{equation}\label{split}
 n\equiv-uv^{-1}\pmod\ell\qquad\text{and}\qquad u+nv\equiv0\pmod k .
\end{equation}
The second condition depends on $n$ only modulo $k$.  Fix $K\ge1$ and put
\[
 \begin{gathered}
 \cK(K)=\{k\le K:\ k\equiv1\ (4)\},\qquad L_K=\lcm\,\cK(K),\\
 h(K)=\sum_{k\in\cK(K)}\frac{\varphi(k)}{k^2}.
 \end{gathered}
\]
If $n\equiv c\pmod{L_K}$, the second condition in \eqref{split} reads
$u+cv\equiv0\pmod k$ and is a condition on $(u,v,k)$ alone; then every $n$ in
the progression $c\pmod{L_K}$ lying in the class $-uv^{-1}\pmod\ell$ is
solvable.

\section{Lattice sums}\label{sec:lattice}

For $k\ge1$, $(c,k)=1$ and $1\le H<z$ let
\[
 \cS(k,c;H,z)=\{(u,v):\ H<u,v\le z,\ (u,v)=1,\ (v,k)=1,\ u\equiv-cv\ (k)\}.
\]
Note that $(u,k)=1$ automatically.  Write $\Lambda=\log(z/H)$.

\begin{lemma}\label{lem:lattice}
There is an absolute constant $K_0$ such that the following holds.  Let $K\ge K_0$, $k\le K$, $(c,k)=1$, $H=K^{10}$ and $z\ge H^2$.  Then
\begin{align}
 \sum_{(u,v)\in\cS(k,c;H,z)}\frac1{uv}&\ \ge\ \frac1{4}\,\frac{\varphi(k)}{k^2}\,\Lambda^2,
 \label{lat-lower}\\
 \sum_{(u,v)\in\cS(k,c;H,z)}\frac1{\varphi(u)\varphi(v)}&\ \le\ C_1\,\frac{\varphi(k)}{k^2}\,\Lambda^2
 \label{lat-upper}
\end{align}
with an absolute constant $C_1$.
\end{lemma}

\begin{proof}
We use two elementary estimates, valid for $1\le H'<z'$, $\Lambda'=\log(z'/H')$,
$(a,k)=1$:
\begin{align}
 \sum_{\substack{H'<u\le z'\\ u\equiv a\,(k)}}\frac1u
 &=\frac{\Lambda'}{k}+O\Big(\frac1{H'}\Big),\label{ap1}\\
 \sum_{\substack{H'<v\le z'\\ (v,k)=1}}\frac1v
 &=\frac{\varphi(k)}{k}\Lambda'+O\Big(\frac{\tau(k)}{H'}\Big).\label{ap2}
\end{align}
For \eqref{ap1}, the terms are $u_0+jk$ with $H'<u_0\le H'+k$, and comparison
with $\frac1k\int dt/t$ gives the sum within $1/H'$ of $\Lambda'/k$ (for a
nonempty sum the upper bound is $1/u_0+\frac1k\log(z'/u_0)$; the lower bound is
$\frac1k\log(z'/(H'+k))\ge\Lambda'/k-1/H'$).
For an empty sum, $z'<u_0\le H'+k$ gives $\Lambda'/k\le1/H'$ directly.
For \eqref{ap2}, M\"obius inversion over $e\mid k$ and
$\sum_{A<m\le B}1/m=\log(B/A)+O(1/A)$ give
$\sum_{e\mid k}\frac{\mu(e)}e\big(\Lambda'+O(e/H')\big)$.

\emph{Lower bound.}  Let $T(H',z')$ denote the sum of $1/(uv)$ over
$H'<u,v\le z'$ with $(v,k)=1$ and $u\equiv-cv\pmod k$ (no coprimality of
$u,v$).  Summing \eqref{ap1} over $v$ and then using \eqref{ap2},
\[
 T(H',z')=\frac{\varphi(k)}{k^2}\Lambda'^2+O\Big(\frac{\Lambda'\tau(k)}{H'}\Big).
\]
valid for $H'\ge1$ and $\Lambda'\ge1$: the additional product error
$O(\tau(k)/H'^2)$ is then absorbed.  M\"obius inversion over $d=(u,v)$ (necessarily $(d,k)=1$, and $d$ is
invertible modulo $k$, so the congruence descends to $u'\equiv-cv'$) gives
\[
 \sum_{\cS}\frac1{uv}=\sum_{\substack{d\le z\\(d,k)=1}}\frac{\mu(d)}{d^2}\,T(H/d,z/d),
\]
and $\log\big((z/d)/(H/d)\big)=\Lambda$ for every $d$.  For $d>H$ we use
only the trivial bound $T(H/d,z/d)\le(1+\log z)^2$ (each one-variable sum
is at most $1+\log z$), which contributes at most
$(1+\log z)^2\sum_{d>H}d^{-2}\ll\Lambda^2/H$ in total, using
$\log z\le2\Lambda$ (from $z\ge H^2$).  Hence
\[
 \sum_{\cS}\frac1{uv}
 =\frac{\varphi(k)}{k^2}\Lambda^2\sum_{\substack{d\le H\\(d,k)=1}}\frac{\mu(d)}{d^2}
 +O\Big(\frac{\Lambda\tau(k)\log z}{H}+\frac{\Lambda^2}{H}\Big).
\]
The sum over $d$ is $\prod_{p\nmid k}(1-p^{-2})+O(1/H)\ge 6/\pi^2-O(1/H)$.
Since $\tau(k)k^2/\varphi(k)\le k\tau(k)^2\le4K^2$ (using
$k/\varphi(k)\le2^{\omega(k)}\le\tau(k)\le2\sqrt k$), the whole error is at most
$CK^2H^{-1}\cdot\frac{\varphi(k)}{k^2}\Lambda^2= CK^{-8}\,\frac{\varphi(k)}{k^2}\Lambda^2$
for an absolute $C$.  This gives \eqref{lat-lower} for $K\ge K_0$, with an
absolute $K_0$; below we only use $K\to\infty$, so we assume $K\ge K_0$ from
now on.

\emph{Upper bound.}  Drop $(u,v)=1$ and use $1/\varphi(m)=\frac1m\sum_{d\mid m}\mu^2(d)/\varphi(d)$.
For $(a,k)=1$, only $d$ with $(d,k)=1$ contribute.  By \eqref{ap1} for
$d\le H$, and by the trivial bound $1+\log z\le1+2\Lambda$ for $d>H$,
\[
\begin{aligned}
 \sum_{\substack{H<u\le z\\ u\equiv a\,(k)}}\frac1{\varphi(u)}
 &=\sum_{\substack{d\le z\\(d,k)=1}}\frac{\mu^2(d)}{d\varphi(d)}
   \sum_{\substack{H/d<u'\le z/d\\ u'\equiv ad^{-1}\,(k)}}\frac1{u'}\\
 &\le\frac{\Lambda}{k}\sum_d\frac{\mu^2(d)}{d\varphi(d)}+\frac1H\sum_{d\le H}\frac{\mu^2(d)}{\varphi(d)}
 +(1+2\Lambda)\sum_{d>H}\frac{1}{d\varphi(d)}\\
 &\le C\Big(\frac{\Lambda}{k}+\frac{\log H}{H}+\frac{\Lambda}{\sqrt H}\Big),
\end{aligned}
\]
because
\[
 \sum_d\frac{\mu^2(d)}{d\varphi(d)}<\infty,\qquad
 \sum_{d\le H}\frac{\mu^2(d)}{\varphi(d)}\ll\log H,\qquad
 \sum_{d>H}\frac1{d\varphi(d)}\ll H^{-1/2}.
\]
Likewise, by \eqref{ap2},
\[
 \sum_{\substack{H<v\le z\\(v,k)=1}}\frac1{\varphi(v)}
 \le C\Big(\frac{\varphi(k)}{k}\Lambda+\frac{\tau(k)\log H}{H}+\frac{\Lambda}{\sqrt H}\Big).
\]
Since $k\le K$, $\tau(k)\le2\sqrt K$, $\varphi(k)/k\ge1/\tau(k)\ge1/(2\sqrt K)$,
$H=K^{10}$ and $\Lambda\ge\frac12\log z\ge10\log K$, the secondary terms in
both brackets are $O(K^{-4})$ times the main terms.  Multiplying proves
\eqref{lat-upper}.
\end{proof}

\begin{lemma}\label{lem:h}
$h(K)\asymp\log K$ for $K\ge5$.
\end{lemma}

\begin{proof}
Writing $\varphi(k)/k=\sum_{d\mid k}\mu(d)/d$,
\[
 h(K)=\sum_{\substack{d\le K\\d\ \mathrm{odd}}}\frac{\mu(d)}{d^2}\sum_{\substack{m\le K/d\\ m\equiv d^{-1}\,(4)}}\frac1m
 =\sum_{\substack{d\le K\\d\ \mathrm{odd}}}\frac{\mu(d)}{d^2}\Big(\frac14\log\frac Kd+O(1)\Big)
\]
and the last expression equals $\frac{2}{\pi^2}\log K+O(1)$,
because $\sum_{d\,\mathrm{odd}}\mu(d)d^{-2}=\prod_{p>2}(1-p^{-2})=8/\pi^2$
and $\sum_d\mu^2(d)d^{-2}(1+\log d)<\infty$.
\end{proof}

\section{The class-mass lemma}\label{sec:mass}

Given $X$, fix the blocks $(x_j,2x_j]$ with $x_j=2^jX^{1/2}$ and
$0\le j<J=\lceil\log X/(2\log2)\rceil$; the last may extend beyond $X$.
Fix $K\ge K_0$, put $H=K^{10}$, and fix $c$ with $(c,L_K)=1$.  For a prime
$\ell\equiv3\pmod 4$ lying in one of these blocks $(x,2x]$, put $z=x^{1/6}$ and
let $\Omega_c(\ell)$ be the set of residues $-uv^{-1}\pmod\ell$ arising
from triples
\begin{equation}\label{triples}
 k\in\cK(K),\qquad (u,v)\in\cS(k,c;H,z),\qquad 4uv\mid k\ell+1 .
\end{equation}
By Lemma~\ref{lem:identity} and \eqref{split} (with $w=(k\ell+1)/(4uv)$),
every $n\equiv c\pmod{L_K}$ with $n\bmod\ell\in\Omega_c(\ell)$ satisfies
\eqref{ES}.  Let $\omega_c(\ell)=|\Omega_c(\ell)|$, and set
$\Omega_c(\ell)=\varnothing$, $\omega_c(\ell)=0$ for primes $\ell\not\equiv3\pmod4$.

\begin{lemma}\label{lem:mass}
There are absolute constants $a,A>0$ and $X_0$ such that the following holds.
Let $X\ge X_0$, $K_0\le K\le(\log X)^4$ (with $K_0$ from Lemma~\ref{lem:lattice}), $H=K^{10}$, $(c,L_K)=1$.  Then, for every
prime $\ell\in(X^{1/2},X]$, $\omega_c(\ell)\le\ell^{1/3}$, and
\begin{equation}\label{mass}
 a(\log X)^2h(K)\ \le\ \sum_{\substack{X^{1/2}<\ell\le X\\ \ell\equiv3\,(4)}}\frac{\omega_c(\ell)}{\ell}\ \le\ A(\log X)^2h(K).
\end{equation}
\end{lemma}

\begin{proof}
\emph{Distinctness.}  Fix $\ell\in(x,2x]$.  Two triples $(k,u,v)$, $(k',u',v')$
in \eqref{triples} with the same residue satisfy $\ell\mid uv'-u'v$ and
$|uv'-u'v|<z^2=x^{1/3}<\ell$, so $uv'=u'v$ and, both pairs being coprime,
$(u,v)=(u',v')$.  Then $k\ell\equiv k'\ell\equiv-1\pmod{4uv}$ forces
$k\equiv k'\pmod{4uv}$, and $4uv>4H^2>K$ gives $k=k'$.  Hence
$\omega_c(\ell)$ equals the number of triples \eqref{triples}, and
$\omega_c(\ell)\le z^2\le\ell^{1/3}$.

\emph{Dyadic blocks.}  Use the fixed blocks above, taking the last block
in full for the upper bound and discarding it for the lower bound if it
extends beyond $X$.  There are $\asymp\log X$ blocks in either case, and
$\log x\asymp\log X$ in each.  Since $K\le(\log X)^4$ and $X$ is large, $z=x^{1/6}\ge H^2=K^{20}$,
so Lemma~\ref{lem:lattice} applies with $\Lambda=\log(z/H)\ge\frac12\log z=\frac1{12}\log x$.
For a triple, the condition $4uv\mid k\ell+1$ is the reduced progression
\begin{equation}\label{prog}
 \ell\equiv-k^{-1}\pmod{4uv},\qquad 4uv\le4x^{1/3},
\end{equation}
which already implies $\ell\equiv3\pmod4$.  Thus
\begin{equation}\label{count}
 \sum_{x<\ell\le2x}\omega_c(\ell)=\sum_{k\in\cK(K)}\ \sum_{(u,v)\in\cS(k,c;H,z)}
 \big(\pi(2x;4uv,-k^{-1})-\pi(x;4uv,-k^{-1})\big).
\end{equation}

\emph{Upper bound.}  By Brun--Titchmarsh, $\pi(2x;q,a)\ll x/(\varphi(q)\log x)$
for $q\le x^{1/2}$, and $\varphi(4uv)\ge2\varphi(u)\varphi(v)$; so by
\eqref{lat-upper}
\[
 \sum_{x<\ell\le2x}\omega_c(\ell)\ll\frac{x}{\log x}\sum_{k}\sum_{\cS}\frac1{\varphi(u)\varphi(v)}
 \ll\frac x{\log x}\,(\log x)^2\,h(K)=x\log x\,h(K).
\]
Dividing by $x$ and summing over the blocks gives the upper half of \eqref{mass}.

\emph{Lower bound.}  Write
$E(x;q,a)=\pi(2x;q,a)-\pi(x;q,a)-(\li(2x)-\li(x))/\varphi(q)$ and
$E^*(q)=\max_{(a,q)=1}|E(x;q,a)|$.  The main term of \eqref{count} is, by
$\li(2x)-\li(x)\ge x/\log(2x)$, $1/\varphi(4uv)\ge1/(4uv)$ and \eqref{lat-lower},
\[
 \ge\frac{x}{4\log 2x}\sum_{k\in\cK(K)}\frac14\frac{\varphi(k)}{k^2}\Lambda^2
 \ \ge\ c_0\,x\log x\,h(K)
\]
for an absolute $c_0>0$.  For the error, group the triples by $q=4uv$: a given
$q$ arises from at most $K\tau(q)$ triples (at most $K$ values of $k$ and
$\tau(q/4)$ ordered factorizations $uv=q/4$).  By Cauchy--Schwarz, the trivial
bound $E^*(q)\ll x/(\varphi(q)\log x)$ (Brun--Titchmarsh), the estimate
$\sum_{q\le Q}\tau(q)^2/\varphi(q)\ll(\log Q)^4$, and the Bombieri--Vinogradov
theorem in the form $\sum_{q\le x^{1/2}(\log x)^{-B}}E^*(q)\ll_R x(\log x)^{-R}$,
where $B=B(R)$ (note $4x^{1/3}$ lies below this level for large $x$),
\[
\begin{aligned}
 \sum_{q\le4x^{1/3}}K\tau(q)E^*(q)
 &\le K\Big(\sum_q\tau(q)^2E^*(q)\Big)^{1/2}\Big(\sum_qE^*(q)\Big)^{1/2}\\
 &\ll K\big(x(\log x)^3\cdot x(\log x)^{-R}\big)^{1/2}
 \le\frac{x}{(\log x)^{3}}
\end{aligned}
\]
on taking $R=19$ and using $K\le(\log X)^4\le 32(\log x)^4$.  Since
$h(K)\ge1$ (the term $k=1$), the error is $o(x\log x\,h(K))$ and
\[
 \sum_{x<\ell\le2x}\frac{\omega_c(\ell)}{\ell}\ \ge\ \frac1{2x}\sum_{x<\ell\le2x}\omega_c(\ell)\ \ge\ \frac{c_0}{3}\log x\,h(K).
\]
Summing over the $\asymp\log X$ full blocks contained in $(X^{1/2},X]$
gives the lower half of \eqref{mass}.
For completeness, the divisor estimate used above follows by bounding its
sum by the Euler product with local factors
$1+\sum_{j\ge1}(j+1)^2/\varphi(p^j)=1+4/p+O(p^{-2})$ and applying Mertens' bound.
The quoted Bombieri--Vinogradov form follows from the standard version with
$\max_{(a,q)=1}\max_{2\le y\le2x}|\pi(y;q,a)-\li(y)/\varphi(q)|$;
the two endpoints cost only a factor $2$.  The bound for $E^*(q)$ includes
the main term, since $\li(2x)-\li(x)\ll x/\log x$.
\end{proof}

The point of Lemma~\ref{lem:mass} is its uniformity in $c$: the only place
$c$ entered was Lemma~\ref{lem:lattice}, whose bounds hold for every reduced
residue.  Without multipliers (only $k=1$) the argument is Lemma 4.1 of \cite{PW}.

\section{The large sieve; proof of Theorem \ref{thm:main}}\label{sec:sieve}

We use Montgomery's large sieve \cite[\S8, p.~561]{Montgomery} in the following
form.  Let $\cN$ be a set of integers in an interval of length $Y$, and let
$\mathcal P$ be a finite set of primes such that, for each $\ell\in\mathcal P$, the
elements of $\cN$ avoid $0\le\omega(\ell)<\ell$ residue classes modulo $\ell$.  Then
for every $Q\ge1$
\begin{equation}\label{LS}
 |\cN|\le\frac{Y+Q^2}{S(Q)},\qquad
 S(Q)=\sum_{\substack{s\le Q\\ s\mid\prod\mathcal P}}\mu^2(s)\prod_{\ell\mid s}g(\ell),
 \qquad g(\ell)=\frac{\omega(\ell)}{\ell-\omega(\ell)} .
\end{equation}
We also use the standard lower bound for $S(Q)$ (as in \cite[\S4]{PW}).  Put
$G=\prod_{\ell\in\mathcal P}(1+g(\ell))=\sum_{s\mid\prod\mathcal P}\mu^2(s)\prod_{\ell\mid s}g(\ell)$.
For any $v>0$, Rankin's trick gives
\[
 \frac{G-S(Q)}{G}\le Q^{-v}\prod_{\ell\in\mathcal P}\frac{1+g(\ell)\ell^{v}}{1+g(\ell)}
 \le Q^{-v}\exp\Big\{\sum_{\ell\in\mathcal P}g(\ell)(\ell^v-1)\Big\}.
\]
If $\mathcal P\subseteq(1,X]$ and $v=1/\log X$, then $\ell^v-1\le\e-1$, and
if moreover $\omega(\ell)\le\ell/2$ then $g(\ell)\le2\omega(\ell)/\ell$; hence
\begin{equation}\label{tail}
 \frac{G-S(Q)}{G}\le\exp\Big\{-\frac{\log Q}{\log X}+2(\e-1)\sum_{\ell\in\mathcal P}\frac{\omega(\ell)}\ell\Big\}.
\end{equation}

\begin{proof}[Proof of Theorem \ref{thm:main}]
Let $N$ be large, fix an absolute $0<\delta<1/4$, and put
$K=\lfloor\delta\log N\rfloor$, $M=L_K$.  Since $L_K\mid\lcm(1,\dots,K)=\e^{\psi(K)}$
and $\psi(K)\le2K$ (Chebyshev), $M\le N^{2\delta}$, so $Y:=N/M\ge N^{1-2\delta}$.
Put
\begin{equation}\label{X}
 X=\exp\Big\{\alpha\Big(\frac{\log N}{\log\log N}\Big)^{1/3}\Big\}
\end{equation}
with a small absolute $\alpha>0$ to be chosen; then $K\le(\log X)^4$ for
large $N$, so Lemma~\ref{lem:mass} applies.

Every prime $p>K$ is coprime to $M$.  Fix a reduced representative
$1\le c\le M$ and let $\cN_c$ be the set of integers $m\in[0,Y]$ for which
$c+Mm$ is an exceptional prime not exceeding $N$.  For each prime
$\ell\in(X^{1/2},X]$, $\ell\equiv3\pmod 4$, we have $\ell\nmid M$ (as
$\ell>K$), and by Lemma~\ref{lem:identity} the number $n=c+Mm$ is solvable
whenever $n\bmod\ell\in\Omega_c(\ell)$; so $m$ avoids the
$\omega_c(\ell)$ classes $(\rho-c)M^{-1}\pmod\ell$, $\rho\in\Omega_c(\ell)$.
Apply \eqref{LS} with $Q=Y^{1/2}$: $|\cN_c|\le2Y/S(Q)$.  By
Lemma~\ref{lem:mass}, $\omega_c(\ell)\le\ell^{1/3}\le\ell/2$ and
$\sum_\ell\omega_c(\ell)/\ell\le A(\log X)^2h(K)\le A'(\log X)^2\log K$
(Lemma~\ref{lem:h}); so \eqref{tail} gives
\[
 \frac{G-S(Q)}{G}\le\exp\Big\{-\frac{(1-2\delta)\log N}{2\log X}+2(\e-1)A'(\log X)^2\log K\Big\}\le\frac12
\]
provided $\alpha$ is so small that $(\log X)^3\log K\le(\log N)/(16(\e-1)A')$;
by \eqref{X} and $\log K\le\log\log N$ this holds for all large $N$.
Therefore, using $\log(1+g)\ge g/2\ge\omega/(2\ell)$ and the lower half of \eqref{mass},
\[
 S(Q)\ge\frac G2\ge\frac12\exp\Big\{\frac12\sum_\ell\frac{\omega_c(\ell)}{\ell}\Big\}
 \ge\frac12\e^{a(\log X)^2h(K)/2}
 \ge\frac12\e^{a'(\log X)^2\log K},
\]
uniformly in $c$.  Summing over the $\varphi(M)\le M$ reduced classes and
adding the primes $p\le K$,
\[
 E_{\mathrm{pr}}(N)\le K+M\cdot\frac{4Y}{\e^{a'(\log X)^2\log K}}
 \ll N\exp\Big\{-a'\alpha^2\,\frac{(\log N)^{2/3}\log(\delta\log N)}{(\log\log N)^{2/3}}\Big\},
\]
which is the theorem.
\end{proof}

\section{All denominators; proof of Theorem \ref{thm:all}}\label{sec:all}

If $p\mid n$ and $4/p=1/x_1+1/x_2+1/x_3$ then
$4/n=\sum_i1/((n/p)x_i)$.  Hence every exceptional $n>1$ has all its prime
factors exceptional, i.e.\ lies in the multiplicative semigroup $\cS$
generated by the exceptional primes, including the empty product $1$.
Put $g(u)=u^{2/3}(\log u)^{1/3}$ and choose $c_0>0$ small enough that
Theorem~\ref{thm:main} gives
$E_{\mathrm{pr}}(t)\le t\exp\{-c_0g(\log t)\}$ for $t\ge t_0\ge\exp(\e)$.  Fix
$0<\eta<c_0$, let $x$ be large, $L=\log x$, and $\delta=\eta g(L)/L\in(0,1/4)$.
By Rankin's trick,
\[
 E(x)\le\sum_{\substack{n\in\cS\\ n\le x}}\Big(\frac xn\Big)^{1-\delta}
 \le x^{1-\delta}\prod_{\substack{p\le x\\ p\ \mathrm{exceptional}}}\big(1-p^{-1+\delta}\big)^{-1}.
\]
The product is $\ll\exp\{\sum_{p\le x\ \mathrm{exc.}}p^{-1+\delta}\}$, since
higher prime powers contribute $O(\sum_p p^{-3/2})=O(1)$ uniformly for $\delta<1/4$.
The endpoint in partial summation is
$E_{\mathrm{pr}}(x)x^{-1+\delta}\le\exp\{-(c_0-\eta)g(L)\}\le1$; hence
\[
 \sum_{\substack{t_0<p\le x\\ p\ \mathrm{exc.}}}p^{-1+\delta}
 \ll1+\int_{t_0}^{x}E_{\mathrm{pr}}(t)\,t^{-2+\delta}\,dt
 \le1+\int_{\log t_0}^{L}\exp\{\delta u-c_0g(u)\}\,du .
\]
Since $g(u)/u$ is decreasing for $u\ge\e$, we have $\delta u\le\eta g(u)$ for
$\log t_0\le u\le L$, so the integral is at most
$\int_{\e}^\infty\exp\{-(c_0-\eta)g(u)\}du<\infty$.  Thus the Euler product is
bounded uniformly in $x$, and $E(x)\ll x^{1-\delta}=x\exp\{-\eta g(\log x)\}$.
\qed

\section{Remarks}

\begin{remark}[general numerators]
Replacing $4$ by a fixed $m\ge4$ (the Erd\H{o}s--Straus--Schinzel problem
$m/n=1/x+1/y+1/z$), the identity of Lemma~\ref{lem:identity} becomes: if
$k\ell\equiv-1\pmod m$, $A=(k\ell+1)/m=uvw$ and $nv\equiv-u\pmod{k\ell}$, then
$m/n=1/(suw)+1/(nsvw)+1/(nuvw)$, where $s=(nv+u)/(k\ell)$.
Take $k\equiv1\pmod m$, $\ell\equiv-1\pmod m$, and replace $4uv$ by $muv$
in the class count.  The lattice lemma is unchanged, while
$\sum_{k\le K,\ k\equiv1\,(m)}\varphi(k)/k^2\asymp_m\log K$
by the proof of Lemma~\ref{lem:h}; all remaining constants may depend on $m$.
The argument above then gives
$E_m(N)\ll_m N\exp\{-c_m(\log N)^{2/3}(\log\log N)^{1/3}\}$; we have not
tracked the dependence on $m$, which \cite{PW} make explicit for Vaughan's
bound.
\end{remark}

\begin{remark}[on the method]
Vaughan's argument is a large sieve with prime moduli $\ell$ and, on average,
$\asymp(\log\ell)^2$ excluded classes per modulus.  Lemma~\ref{lem:mass}
raises the class mass from $(\log X)^2$ to $(\log X)^2h(K)$ in its stated
range.  Extrapolating formally to $\log K\asymp\log X$ suggests a cubic
class mass, but this is not established here.  Nor do the prime-average
solution counts of \cite{ET} establish a typical order $\asymp(\log n)^3$.
The present note uses a single progression modulo $L_K\le N^{1/2}$ to
absorb the $k$-part, which produces the $\log\log$ factor.  A cubic class
mass with comparable sieve and tail control would give an exponent $3/4$;
obtaining such control is not addressed here.
\end{remark}

\begin{thebibliography}{9}
\bibitem{ET} C.~Elsholtz and T.~Tao, Counting the number of solutions to the
Erd\H{o}s--Straus equation on unit fractions, J.\ Aust.\ Math.\ Soc.\ 94 (2013),
50--105.
\bibitem{Montgomery} H.~L.~Montgomery, The analytic principle of the large
sieve, Bull.\ Amer.\ Math.\ Soc.\ 84 (1978), 547--567.
\bibitem{Mordell} L.~J.~Mordell, Diophantine Equations, Academic Press, 1969,
Chapter~30.
\bibitem{PW} C.~Pomerance and A.~Weingartner, Exceptions to the
Erd\H{o}s--Straus--Schinzel conjecture, arXiv:2511.16817 (2025).
\bibitem{Salez} S.~E.~Salez, The Erd\H{o}s--Straus conjecture: new modular
equations and checking up to $N=10^{17}$, arXiv:1406.6307 (2014).
\bibitem{Vaughan} R.~C.~Vaughan, On a problem of Erd\H{o}s, Straus and
Schinzel, Mathematika 17 (1970), 193--198.
\end{thebibliography}

\end{document}
